\section{Related Work}\label{sec:related}
\subsection{Fraud detection and class imbalance}
Statistical and data-mining approaches to financial fraud are well surveyed~\cite{bolton,ngai}. The defining difficulty is extreme class imbalance, which makes accuracy and ROC-AUC misleading and motivates precision--recall metrics and cost-aware evaluation~\cite{haixiang}. The public credit-card benchmark of Dal~Pozzolo et~al.~\cite{dalpozzolo}, which has a fraud prevalence of about 0.17\%, is widely used for this purpose, and gradient-boosted trees such as XGBoost~\cite{xgboost} are strong, fast baselines on tabular data of this type. Studies reporting quantum gains on such data rarely control for model capacity; this gap motivates the capacity-matched controls used here.

\subsection{Quantum machine learning}
Variational quantum classifiers and quantum-kernel methods~\cite{havlicek,schuld} are the main QML families applied to tabular classification. A known obstacle is trainability: Cerezo et~al.~\cite{cerezo} show that cost functions built from \emph{global} observables can exhibit barren plateaus even in shallow circuits, whereas \emph{local} observables are expected to be more trainable. Kyriienko and Magnusson~\cite{kyriienko} propose an unsupervised quantum-kernel anomaly detector evaluated on the same credit-card dataset. Entanglement structure can also be learned from data, e.g.\ with Chow--Liu dependence trees~\cite{chowliu}. These results motivate the observable-locality and structured-circuit hypotheses tested in Layer~1, and they do not by themselves establish an advantage over strong classical models.

\subsection{Post-quantum cryptography}
Shor's algorithm~\cite{shor} threatens RSA and elliptic-curve cryptography, while Grover's algorithm~\cite{grover} reduces generic symmetric-key search cost quadratically; the ``harvest now, decrypt later'' concern motivates early migration~\cite{mosca,nistir8105}. NIST has standardized ML-KEM~\cite{nist203} (from CRYSTALS-Kyber~\cite{kyber}), ML-DSA~\cite{nist204} (from CRYSTALS-Dilithium~\cite{dilithium}) and SLH-DSA~\cite{nist205} (from SPHINCS+~\cite{sphincs}). Their primitives differ in size and speed: lattice schemes are fast with moderate sizes, whereas hash-based signatures are conservative in assumptions but large and slow to sign. Reported speeds depend strongly on the implementation, which motivates the reference-versus-compiled comparison in this work.

\subsection{Quantum key distribution and NISQ hardware}
BB84~\cite{bb84} underlies most QKD; security proofs~\cite{shorpreskill} and reviews~\cite{gisin} establish that, with ideal devices, an intercept--resend attacker induces about 25\% error in the sifted key. Noisy intermediate-scale quantum (NISQ) devices~\cite{preskill} have gate, readout and decoherence errors, so circuit-level execution on such devices measures hardware error rather than provides secure communication. Qiskit~\cite{qiskit} is the software stack used here.

\subsection{Research gap}
Prior work tends to treat fraud detection, post-quantum migration and QKD separately, and quantum-ML claims are seldom tested against capacity-matched classical controls or with pre-specified protocols. This study does not claim priority for any single component. Its contribution is methodological: a reproducible, layered banking-security experiment that (i) tests a quantum-model design choice with paired statistics and a classical control, (ii) separates implementation-backend cost from primitive cost in post-quantum cryptography, (iii) cross-validates BB84 at protocol and circuit level, and (iv) quantifies simulator-to-hardware deviation on IBM Quantum devices.
